\documentclass[../../main.tex]{subfiles}
\begin{document}
\section{同伦不变性}\label{section:ch17-homotopy-invariance}

本节我们证明一个重要的推广：de Rham 群实际上是同伦不变量，这特别意味着它们是拓扑不变量。

\begin{theorem}[同伦映射诱导相同的上同调映射]\label{theorem:ch17-homotopy-invariance}
设 $F,G:M\to N$ 是（带边或不带边）光滑流形之间的光滑映射。若 $F$ 与 $G$ 光滑同伦，则它们在 de Rham 上同调上诱导相同的映射：
\begin{align*}
F^*=G^*:H^p_{\mathrm{dR}}(N)\to H^p_{\mathrm{dR}}(M)
\end{align*}
对所有 $p$ 成立。
\end{theorem}

\begin{proposition}[同伦算子]\label{proposition:ch17-homotopy-operator}
设 $M,N$ 是光滑流形，$H:M\times I\to N$ 是光滑同伦，连接 $F=H_0$ 与 $G=H_1$。则存在一个线性映射（同伦算子）$\mathcal{h}:\Omega^p(N)\to\Omega^{p-1}(M)$，使得对每个闭形式 $\omega$ 有
\begin{align}
G^*\omega-F^*\omega=\mathrm{d}(\mathcal{h}\omega)+\mathcal{h}(\mathrm{d}\omega).
\end{align}
于是，当 $\omega$ 闭时，$G^*\omega-F^*\omega$ 是恰当的。
\end{proposition}

\begin{figure}[ht]
\centering
\includegraphics[width=0.6\textwidth]{fig-17-1.png}
\caption{同伦不变性}\label{figure:ch17-homotopy-invariance}
\end{figure}

\begin{proof}
把 $H:M\times I\to N$ 写成 $H(p,t)$。令 $i_0,i_1:M\to M\times I$ 分别为 $i_0(p)=(p,0)$，$i_1(p)=(p,1)$。在 $M\times I$ 上引入坐标 $(x^1,\dots,x^n,t)$，把任意形式分解为不含 $\mathrm{d}t$ 的部分与含一个 $\mathrm{d}t$ 的部分。利用在 $\mathrm{d}t$ 方向上的“纤维积分”（关于 $t$ 从 $0$ 到 $1$ 积分）定义 $\mathcal{h}$。直接计算表明 $\mathrm{d}\mathcal{h}+\mathcal{h}\mathrm{d}=i_1^*-i_0^*$。由于 $F=H\circ i_0$，$G=H\circ i_1$，这就给出所要的恒等式。
\end{proof}

\begin{corollary}[Poincaré 引理的上同调形式]\label{corollary:ch17-poincare-lemma-cohomology}
设 $U\subseteq\mathbb{R}^n$ 是一个星形开集。则 $H^p_{\mathrm{dR}}(U)=0$ 对所有 $p\geqslant1$ 成立，而 $H^0_{\mathrm{dR}}(U)=\mathbb{R}$。
\end{corollary}

\begin{proof}
恒等映射 $\operatorname{Id}:U\to U$ 与映到一点的常值映射同伦（利用从中心点出发的径向收缩）。由同伦不变性，$\operatorname{Id}^*$ 与常值映射诱导的上同调映射相同。而常值映射在 $p\geqslant1$ 次上诱导零映射，故 $H^p_{\mathrm{dR}}(U)=0$。
\end{proof}

\begin{corollary}
若两个流形 $M,N$ 光滑同伦等价，则 $H^p_{\mathrm{dR}}(M)\cong H^p_{\mathrm{dR}}(N)$ 对所有 $p$ 成立。特别地，微分同胚的流形有相同的 de Rham 群。
\end{corollary}
\end{document}
